\documentclass[11pt,a4paper]{article}
\usepackage[utf8]{inputenc}
\usepackage[T1]{fontenc}
\usepackage{geometry}
\geometry{margin=1in}
\usepackage{hyperref}
\usepackage{enumitem}
\usepackage{booktabs}
\usepackage{longtable}
\usepackage{xcolor}
\usepackage{titlesec}
\usepackage{amsmath,amssymb}

\titleformat{\section}{\Large\bfseries\color{blue!70!black}}{}{0em}{}
\titleformat{\subsection}{\large\bfseries\color{blue!50!black}}{}{0em}{}
\titleformat{\subsubsection}{\normalsize\bfseries\color{blue!30!black}}{}{0em}{}

\title{\textbf{Replication Plan}\\[0.5em]
\large NukeLM: Pre-Trained and Fine-Tuned Language Models for the Nuclear and Energy Domains\\[0.3em]
\normalsize OSTI ID: 1861801}
\author{Auto-generated Replication Blueprint}
\date{\today}

\begin{document}
\maketitle
\tableofcontents
\newpage

%---------------------------------------------------------------------
\section{Paper Overview}
%---------------------------------------------------------------------
NukeLM presents domain-specific language models for the nuclear science and energy domains. Starting from general-purpose pre-trained models (RoBERTa base/large, SciBERT), the authors perform continued pre-training (domain-adaptive pre-training, DAPT) on a large corpus of nuclear/energy-domain text (OSTI abstracts, NRC documents, DOE reports). The models are then fine-tuned on downstream NLP tasks: named entity recognition (NER), relation extraction (RE), and text classification in the nuclear domain. The paper demonstrates that domain-specific pre-training significantly improves performance on these tasks.

%---------------------------------------------------------------------
\section{Detailed Workflow}
%---------------------------------------------------------------------

\subsection{Phase 1: Corpus Collection and Pre-processing}
\begin{enumerate}[label=\textbf{1.\arabic*}]
  \item \textbf{Collect pre-training corpus.}
    \begin{itemize}
      \item OSTI.gov abstracts: query the OSTI API for nuclear science and energy technology abstracts.
      \item NRC ADAMS documents (publicly available regulatory documents).
      \item DOE technical reports from SciTech Connect.
      \item Total corpus size: $\sim$millions of documents / several GB of text.
    \end{itemize}
  \item \textbf{Pre-process text.}
    \begin{itemize}
      \item Clean HTML/XML artifacts, normalize whitespace, remove duplicates.
      \item Sentence tokenization and subword tokenization (matching RoBERTa/SciBERT tokenizer).
    \end{itemize}
\end{enumerate}

\subsection{Phase 2: Domain-Adaptive Pre-Training}
\begin{enumerate}[label=\textbf{2.\arabic*}]
  \item \textbf{Load base models.}
    \begin{itemize}
      \item RoBERTa-base (125M params), RoBERTa-large (355M params), SciBERT.
      \item Use HuggingFace Transformers \texttt{AutoModelForMaskedLM}.
    \end{itemize}
  \item \textbf{Continued pre-training with MLM objective.}
    \begin{itemize}
      \item Masked language modeling on the nuclear domain corpus.
      \item Training parameters: batch size, learning rate, warm-up steps, total steps (as specified in paper).
      \item Use distributed training (multi-GPU) for efficiency.
    \end{itemize}
  \item \textbf{Save domain-adapted model checkpoints} (NukeLM-base, NukeLM-large).
\end{enumerate}

\subsection{Phase 3: Downstream Task Fine-Tuning}
\begin{enumerate}[label=\textbf{3.\arabic*}]
  \item \textbf{Named Entity Recognition (NER).}
    \begin{itemize}
      \item Prepare NER training data with nuclear domain entity tags (e.g., material, isotope, facility).
      \item Add a token classification head on top of NukeLM.
      \item Train and evaluate (precision, recall, F1 per entity type).
    \end{itemize}
  \item \textbf{Relation Extraction (RE).}
    \begin{itemize}
      \item Prepare annotated relation data.
      \item Fine-tune with entity marker strategy or similar.
      \item Evaluate F1 score.
    \end{itemize}
  \item \textbf{Text Classification.}
    \begin{itemize}
      \item Classify documents into nuclear sub-domains.
      \item Add a sequence classification head; fine-tune and evaluate accuracy/F1.
    \end{itemize}
\end{enumerate}

\subsection{Phase 4: Comparison and Validation}
\begin{enumerate}[label=\textbf{4.\arabic*}]
  \item Compare NukeLM vs.\ base RoBERTa, SciBERT, and BioBERT on all downstream tasks.
  \item Reproduce the paper's tables of F1 scores and accuracy.
  \item Statistical significance tests (bootstrap or paired t-test).
\end{enumerate}

%---------------------------------------------------------------------
\section{Datasets}
%---------------------------------------------------------------------

\begin{longtable}{p{4.5cm} p{3.5cm} p{6cm}}
\toprule
\textbf{Dataset / Source} & \textbf{Type} & \textbf{Location / Notes} \\
\midrule
\endfirsthead
\toprule
\textbf{Dataset / Source} & \textbf{Type} & \textbf{Location / Notes} \\
\midrule
\endhead
OSTI.gov abstracts & Pre-training corpus & \url{https://www.osti.gov/api/v1/} (OSTI API) \\
\midrule
NRC ADAMS documents & Pre-training corpus & \url{https://www.nrc.gov/reading-rm/adams.html} \\
\midrule
DOE SciTech Connect & Pre-training corpus & \url{https://www.osti.gov/scitech/} \\
\midrule
Nuclear NER dataset & Fine-tuning (labeled) & Described in paper; may require manual annotation or author-provided data \\
\midrule
Nuclear RE dataset & Fine-tuning (labeled) & Described in paper; check supplementary materials \\
\midrule
Nuclear text classification & Fine-tuning (labeled) & Described in paper; OSTI category labels \\
\bottomrule
\end{longtable}

%---------------------------------------------------------------------
\section{Models, Tools, and Software}
%---------------------------------------------------------------------

\begin{longtable}{p{4.5cm} p{2.5cm} p{6.5cm}}
\toprule
\textbf{Tool / Library} & \textbf{Version} & \textbf{Purpose / URL} \\
\midrule
\endfirsthead
\toprule
\textbf{Tool / Library} & \textbf{Version} & \textbf{Purpose / URL} \\
\midrule
\endhead
HuggingFace Transformers & $\geq 4.15$ & Model loading, training, fine-tuning \newline \url{https://huggingface.co/docs/transformers/} \\
\midrule
HuggingFace Datasets & latest & Dataset loading and preprocessing \\
\midrule
PyTorch & $\geq 1.10$ & Deep learning backend \\
\midrule
RoBERTa (base/large) & pre-trained & Base models from HuggingFace Hub \\
\midrule
SciBERT & pre-trained & Scientific domain base model \newline \url{https://huggingface.co/allenai/scibert_scivocab_uncased} \\
\midrule
seqeval & latest & NER evaluation metrics \\
\midrule
scikit-learn & $\geq 0.24$ & Classification metrics \\
\midrule
Weights \& Biases / TensorBoard & latest & Training monitoring (optional) \\
\bottomrule
\end{longtable}

\subsection{Hardware Requirements}
\begin{itemize}
  \item \textbf{Pre-training}: Multiple NVIDIA A100 or V100 GPUs (32\,GB+); estimated 1--7 days.
  \item \textbf{Fine-tuning}: Single GPU with $\geq$ 16\,GB VRAM.
  \item \textbf{Storage}: $\geq$ 100\,GB for corpus and model checkpoints.
\end{itemize}

\end{document}
